The ridge fits select the regularisation $\alpha$ per voxel and outer fold from a grid whose
lowest value is $10^{-6}$. Following the amendment of 2026-10-03 (Section~\ref{sec:inference}),
the selection of the lowest value at a reliable voxel does not abort the run. It is logged
and, after the run, classified on the mean inner-fold MSE curve of that voxel in that outer
fold. The curve is extended by the diagnostic values $10^{-7}$, $10^{-8}$ and $10^{-9}$, and
the relative gain is $(\mathrm{MSE}(10^{-6}) - \mathrm{MSE}_{\min}) / \mathrm{MSE}(10^{-6})$,
with $\mathrm{MSE}_{\min}$ the minimum over the extended grid. A relative gain below $10^{-3}$
is a plateau and the selection is valid; otherwise the cell would be excluded from every
result. Smaller values of $\alpha$ are not used, because there the solution is dominated by
eigenvalues of the dual kernel at rounding level ($|\lambda| \leq 1.9 \times 10^{-12}$ in the
diagnosed cases). The gate decides whether a cell is kept and never changes a score.

The first cell with such selections was a condition of the pilot (untrained network with
calibrated BN statistics, one seed, $g = 1$, sub-01). There the standardised training design had full column rank
and no dead features (step report in the repository); the pilot diagnosis used a wider
diagnostic grid than the gate. Table~\ref{tab:lowedge} lists every analysis with such
selections. Steps~7 to~9 involve no ridge fits.

\begin{table}[h]
\centering
\footnotesize
\setlength{\tabcolsep}{3pt}
\begin{tabular}{@{}p{0.40\textwidth}rrrr@{}}
\toprule
Analysis & Hits & Plateau & Excluded & Max.\ gain \\
\midrule
Pilot diagnosis (sub-01, $g = 1$) & \LowEdgePilotHits & \LowEdgePilotPlateau & \LowEdgePilotExcluded & \LowEdgePilotMaxGain \\
RSA versus ridge: collapse, random folds (sub-01) & \LowEdgeEtwoHits & \LowEdgeEtwoPlateau & \LowEdgeEtwoExcluded & \LowEdgeEtwoMaxGain \\
RSA versus ridge: run-grouped folds (sub-01) & \LowEdgeEthreeHits & \LowEdgeEthreePlateau & \LowEdgeEthreeExcluded & \LowEdgeEthreeMaxGain \\
RSA versus ridge: encoding-RSA (sub-01; repeats the hits above) & \LowEdgeEfourHits & \LowEdgeEfourPlateau & \LowEdgeEfourExcluded & \LowEdgeEfourMaxGain \\
Other subjects, $g = 1$ (sub-02, sub-03) & \LowEdgeSubjectsHits & \LowEdgeSubjectsPlateau & \LowEdgeSubjectsExcluded & \LowEdgeSubjectsMaxGain \\
Diagnosis of sub-03 (controls, $g \in \{1, 4, 8\}$) & \LowEdgeSubThreeHits & \LowEdgeSubThreePlateau & \LowEdgeSubThreeExcluded & \LowEdgeSubThreeMaxGain \\
Ridge across resolutions (sub-01, sub-02, $g = 1$) & \LowEdgeRidgeHits & \LowEdgeRidgePlateau & \LowEdgeRidgeExcluded & \LowEdgeRidgeMaxGain \\
Pre-specified primary analysis (3 subjects) & \LowEdgePrimaryHits & \LowEdgePrimaryPlateau & \LowEdgePrimaryExcluded & \LowEdgePrimaryMaxGain \\
\bottomrule
\end{tabular}
\caption{Selections of the lowest regularisation value at reliable voxels, classified by the
plateau gate. In the pilot diagnosis the plateau criterion was applied to the relative gain
of the wider diagnostic grid. All hits were plateaus, and no cell was excluded.}
\label{tab:lowedge}
\end{table}
